\chapter{从曲面坐标到切空间}
\label{chap:math-manifolds}

在平面上，位置可以由一对实数表示。球面上却没有一对处处有效的经纬度：
经度在极点失去意义，但球面本身没有裂口。因此需要区分一个点与给它编号的坐标。
我们先说明“局部区域”和“坐标变换”是什么意思，再把求导运算搬到曲面上。

\section{开集和连续性先于坐标图}
\label{sec:math-manifold-topology}

\begin{definition}[拓扑、邻域与闭集]
集合 \(M\) 上的拓扑 \(\mathcal T\) 是一族子集，含
\(\varnothing,M\)，对任意并和有限交封闭。
\(\mathcal T\) 中的集合称为开集；含点 \(p\) 的开集称为 \(p\) 的开邻域；
补集为开集的集合称为闭集。
\end{definition}

\(\mathbb R^n\) 的通常拓扑取所有这样的集合：集合中每个点都能找到一个
完全包含于集合的小开球。于是开区间是 \(\mathbb R\) 的开集，
\([0,1]\) 是闭集，\([0,1)\) 两者都不是。
这个定义只告诉我们哪些点可在局部一起讨论，尚未赋予距离。

\begin{definition}[连续映射与同胚]
拓扑空间间的映射 \(F:M\to N\) 连续，若 \(N\) 中每个开集 \(V\)
的原像 \(F^{-1}(V)\) 都在 \(M\) 中开。
双射 \(F\) 若它和 \(F^{-1}\) 都连续，称为同胚。
\end{definition}

在 Euclid 空间，这与高等数学的 \(\varepsilon\)–\(\delta\) 连续性等价。
例如 \(\mathbb R\) 上的 \(x\mapsto x^3\) 是同胚；
\(x\mapsto x^2\) 在全体实数上不是，因为它不单射。
坐标图要求同胚，原因是反向也必须能从坐标唯一、连续地找回原点。

\begin{definition}[两个分离性条件]
若不同点 \(p,q\) 总有互不相交的开邻域，称 \(M\) 为 Hausdorff 空间。
若存在可数个开集 \(B_1,B_2,\ldots\)，使任意开集都是其中若干项的并，
称 \(M\) 第二可数；\(\{B_j\}\) 称为可数基。
\end{definition}

Hausdorff 条件保证极限唯一：若同一序列同时收敛到 \(p\ne q\)，
最终项必须同时落入两个不相交邻域，矛盾。
第二可数也不是空洞的限制。不可数集合若取离散拓扑，
每个单点都开，任何基必须含全部单点，因而不第二可数。
普通 \(\mathbb R^n\) 却有可数基：中心坐标和半径均取有理数的开球。
若 \(M\subset\mathbb R^N\) 已是几何子集，就取子空间拓扑：
\(U\subset M\) 开，当且仅当存在 \(\mathbb R^N\) 的开集 \(O\)
使 \(U=M\cap O\)。这使球面上“删去一个点后的区域”
有确切的开集意义。

\section{两张图怎样描述同一个球面}

\begin{definition}[坐标图与光滑相容]
\(n\) 维坐标图是 \(M\) 的开集 \(U\) 与同胚
\(\varphi:U\to\varphi(U)\subseteq\mathbb R^n\)，其中右边为开集。
两图 \((U,\varphi),(V,\psi)\) 在重叠处光滑相容，若
\[
\psi\circ\varphi^{-1}:
\varphi(U\cap V)\longrightarrow\psi(U\cap V)
\]
及其逆映射都具有任意阶连续偏导数。
覆盖 \(M\) 且两两相容的图族称为光滑图册。
\end{definition}

要求过渡映射光滑比仅要求每张图连续更强。
在实数轴上，\(\varphi(x)=x\) 与
\(\psi(x)=\sqrt[3]{x}\) 都是到实数轴的同胚，
但 \(\psi\circ\varphi^{-1}(u)=\sqrt[3]{u}\)
在零点不可微。这两张图不能同时放进同一光滑图册。
因此“有局部坐标”与“坐标间可无歧义地求导”是两层不同要求。

取单位球面 \(S^2=\{(x,y,z):x^2+y^2+z^2=1\}\)。
删去北极 \(N=(0,0,1)\) 后，沿北极连线投影到赤道平面，得到
\[
\varphi_N(x,y,z)=\left(\frac{x}{1-z},\frac{y}{1-z}\right).
\]
删去南极 \(S=(0,0,-1)\) 后同样得到
\(\varphi_S(x,y,z)=(x/(1+z),y/(1+z))\)。
两张图都覆盖除相应极点外的整个球面。在重叠区域写
\(u=\varphi_N(p)\)，令 \(r^2=|u|^2\)。由球面方程
\(x^2+y^2=1-z^2\) 得
\(r^2=(1+z)/(1-z)\)，所以
\[
\varphi_N^{-1}(u)=
\left(\frac{2u_1}{1+r^2},
\frac{2u_2}{1+r^2},
\frac{r^2-1}{1+r^2}\right).
\]
分母处处非零；这个显式逆映射与 \(\varphi_N\) 都连续，
因而它确实是坐标图，而不仅是一个看起来合理的投影公式。
在两图重叠处 \(u\ne0\)，代入南极图便得到
\[
(\varphi_S\circ\varphi_N^{-1})(u)=\frac{u}{|u|^2}.
\]
它在 \(u\ne0\) 上光滑，且再作用一次便回到 \(u\)。
球面上的点没有因更换坐标而改变；改变的只是它在两张图中的数对。

\begin{definition}[光滑流形与光滑映射]
Hausdorff、第二可数且带有光滑图册的空间称光滑流形。
两个光滑流形间的映射 \(F:M\to N\) 称光滑，
若在任意相容坐标图中，其局部表达
\(\psi\circ F\circ\varphi^{-1}\) 为通常意义下的光滑映射。
\end{definition}

这个定义不要求流形预先嵌入 Euclid 空间。
球面借助嵌入帮助想象，但时空只需一组相容局部坐标。
对光滑性必须同时查看输入和输出的图；否则一个公式仅说明坐标分量，
尚未说明它定义了哪两个流形之间的映射。

\section{曲线速度与点处求导为何是同一对象}
\label{sec:math-tangent-derivations}

球面上的一条曲线 \(\gamma:(-\epsilon,\epsilon)\to S^2\) 有速度，
但把 \(\dot\gamma(0)\) 看作三维向量只适用于已经嵌入
\(\mathbb R^3\) 的曲面。一个不依赖嵌入的办法是记录曲线经过点
\(p=\gamma(0)\) 时，怎样改变所有光滑函数的值。

\begin{definition}[曲线切向量]
经过 \(p\in M\) 的两条可微曲线 \(\gamma_1,\gamma_2\) 等价，
若在一张经过 \(p\) 的坐标图 \(\varphi\) 中有
\((\varphi\circ\gamma_1)'(0)=(\varphi\circ\gamma_2)'(0)\)。
等价类称为 \(p\) 处的切向量，其集合记为 \(T_pM\)。
\end{definition}

这里的“一张”不会导致坐标依赖。若换到图 \(\psi\)，链式法则给
\[
(\psi\circ\gamma)'(0)
=D(\psi\circ\varphi^{-1})_{\varphi(p)}
(\varphi\circ\gamma)'(0).
\]
两曲线在 \(\varphi\) 中速度相同，则在 \(\psi\) 中也相同。
这个等式同时告诉我们：切向量本身不变，分量被坐标变换的 Jacobian 相乘。

\begin{definition}[点处导子]
点 \(p\) 处的导子是一个把 \(p\) 附近光滑实值函数送到实数的线性映射
\(v\)，满足 Leibniz 法则
\[
v(fg)=f(p)v(g)+g(p)v(f).
\]
只在 \(p\) 附近相同的函数视为同一个输入；这种局部函数称为 \(p\) 处的芽。
\end{definition}

曲线 \(\gamma\) 给出导子
\(v_\gamma(f)=\left.\dd f(\gamma(t))/\dd t\right|_{t=0}\)。
乘积求导法则说明它确实满足上面的定义。反过来，导子并不会包含比曲线速度
更多的信息。

\begin{theorem}[曲线定义与导子定义的对应]
在 \(n\) 维光滑流形 \(M\) 上，每个 \(p\) 处的导子都由唯一的曲线切向量给出。
给定坐标 \(x=(x^1,\ldots,x^n)\)，该导子可唯一写成
\[
v=\sum_{i=1}^n v(x^i)\,\partial_i\big|_p,\qquad
\partial_i\big|_p(f)=\partial_i(f\circ x^{-1})(x(p)).
\]
\end{theorem}
\begin{proof}
先由 Leibniz 法则作用于常数函数 \(1=1\cdot1\) 得 \(v(1)=0\)，
故所有常数的导数为零。把局部坐标中的 \(f\) 限制在 \(x(p)\) 附近的小凸球，
沿连接 \(x(p)\) 与 \(x\) 的直线使用微积分基本定理：
\[
f(x)-f(x(p))
=\sum_i(x^i-x^i(p))f_i(x),\qquad
f_i(x)=\int_0^1\partial_i f(x(p)+t(x-x(p)))\,\dd t.
\]
这里 \(f_i\) 光滑且 \(f_i(x(p))=\partial_i f(x(p))\)。
对等式作用 \(v\)，常数项为零，Leibniz 法则使
\((x^i-x^i(p))v(f_i)\) 在 \(p\) 处消失，得到
\[
v(f)=\sum_i v(x^i)\partial_i f(x(p)).
\]
因此数 \(v^i=v(x^i)\) 决定全部作用。取
\(\gamma(t)=x^{-1}(x(p)+t(v^1,\ldots,v^n))\)，小 \(t\) 时该曲线有定义，
链式法则给 \(v_\gamma(f)=v(f)\)。唯一性由各坐标函数的值确定。
\end{proof}

把球面当作 \(\mathbb R^3\) 的子集时，约束
\(x^2+y^2+z^2=1\) 沿曲线求导给
\(p\cdot\dot\gamma(0)=0\)。
反过来，每个满足 \(p\cdot v=0\) 的三维向量都可作为球面上一条曲线
的初速度，例如把 \(p+tv\) 归一化。
因此 \(T_pS^2=\{v\in\mathbb R^3:p\cdot v=0\}\)。
这算出了切空间，也验证抽象定义与熟悉的切平面一致。

\begin{proposition}[正则约束的切空间]
设 \(F:\mathbb R^N\to\mathbb R^m\) 光滑，
\(M=F^{-1}(0)\)，并且 \(DF_p\) 在每个 \(p\in M\)
的秩为 \(m\)。则 \(M\) 在每个点附近是 \((N-m)\) 维
光滑流形，且
\[
T_pM=\ker DF_p.
\]
\end{proposition}
\begin{proof}
由秩条件，可重新排列环境坐标，使 \(DF_p\) 对最后 \(m\)
个坐标的偏导方阵可逆。多变量微积分的隐函数定理遂把这些坐标
局部写成前 \(N-m\) 个坐标的光滑函数，这给出一张局部图。
若 \(\gamma(t)\subset M\) 且 \(\gamma(0)=p\)，
对 \(F(\gamma(t))=0\) 求导，得
\(DF_p\dot\gamma(0)=0\)，故 \(T_pM\subseteq\ker DF_p\)。
局部图的 \(N-m\) 个坐标速度线性独立并全在该核内；
由秩—零度定理，核维数也是 \(N-m\)，故两者相等。
\end{proof}

球面取 \(F(x,y,z)=x^2+y^2+z^2-1\)，
\(DF_pv=2p\cdot v\)，于是回到上面的切平面。
秩条件不能省略：锥面
\(x^2+y^2-z^2=0\) 在顶点的 \(DF_0=0\)，
形式上的核是整个 \(\mathbb R^3\)，却不存在一个二维光滑切平面。
这就是约束方程虽可写下，几何对象却可能在特殊点变坏的具体例子。

\section{映射怎样推动速度}

若 \(F:M\to N\) 光滑，曲线 \(\gamma\) 经 \(F\) 变成
\(F\circ\gamma\)。由此定义点 \(p\) 处的微分
\[
dF_p:T_pM\to T_{F(p)}N,\qquad
dF_p(v)(h)=v(h\circ F).
\]
它是线性的，且因函数复合满足
\(d(G\circ F)_p=dG_{F(p)}\circ dF_p\)。
在坐标中，\(dF_p\) 正是熟悉的 Jacobian 矩阵；
写清定义域与值域能防止把两个不同点的速度直接相加。

\section{对偶向量与张量}

线性代数中，行向量作用于列向量得到一个数。流形上也一样：
\begin{definition}[余切向量与张量]
\(T_pM\) 上的线性泛函组成余切空间 \(T_p^*M\)。
函数 \(f\) 的微分 \(df_p\) 由 \(df_p(v)=v(f)\) 定义。
一个 \((r,s)\) 型张量是在同一点 \(p\) 处，对 \(r\) 个余切向量和
\(s\) 个切向量分别线性的多重映射，值为实数。
\end{definition}

坐标函数的微分 \(dx^i\) 与 \(\partial_j\) 对偶，因为
\(dx^i(\partial_j)=\delta^i{}_j\)。
若新坐标为 \(y^a=y^a(x)\)，上面的链式法则给
\[
\frac{\partial}{\partial y^a}
=\frac{\partial x^i}{\partial y^a}\frac{\partial}{\partial x^i},
\qquad
dy^a=\frac{\partial y^a}{\partial x^i}dx^i .
\]
两种 Jacobian 互逆，故对偶配对不随坐标改变。
在平面极坐标区域 \(r>0\)，直接微分
\(x=r\cos\theta,y=r\sin\theta\)，得
\[
\partial_r=\cos\theta\,\partial_x+\sin\theta\,\partial_y,
\qquad
\partial_\theta=-r\sin\theta\,\partial_x
+r\cos\theta\,\partial_y.
\]
反过来，由 \(r=\sqrt{x^2+y^2}\) 与局部角函数得到
\[
dr=\frac{x\,dx+y\,dy}{r},
\qquad
d\theta=\frac{-y\,dx+x\,dy}{r^2}.
\]
于是 \(dr(\partial_\theta)=0\)、
\(d\theta(\partial_\theta)=1\)，确实是一对对偶基。
但角度 \(\theta\) 不能在整个穿孔平面选成连续单值函数；
右侧的公式却拼成一个处处光滑的全局一形式。
下一章将证明它虽处处局部等于角度的微分，却不是全局函数的微分。
例如点处的双线性型 \(b_p(v,w)\) 是 \((0,2)\) 型张量；记
\((\alpha\otimes\beta)(v,w)=\alpha(v)\beta(w)\)，则它可写成
\(b=b_{ij}dx^i\otimes dx^j\)，这时才出现指标。
因此“张量”首先是多重线性映射，带指标的数组只是选基后的坐标记录。

\begin{exercise}
在圆柱面 \(x^2+y^2=1\) 的局部坐标 \((\theta,z)\) 中，
求 \(\partial_\theta,\partial_z\) 作为 \(\mathbb R^3\) 向量的表达，
并验证它们满足正则约束切空间的判据。
解释为何 \(\theta\) 仍不能作为覆盖整个圆柱面的单张实值坐标。
\end{exercise}

\section{反对称张量、微分形式与积分}
\label{sec:math-forms-from-tensors}

计算沿曲线的积分时，线元 \(\dd x^i\) 会与函数系数一起变换，
使积分与坐标选择无关。面积积分需要同样的机制，但还要求交换两个方向时
方向反号。这把我们带到反对称张量。

\begin{definition}[微分形式与楔积]
\(p\) 处的 \(k\)-形式是 \(T_pM\) 上的交替 \(k\) 重线性映射：
交换任意两个输入会使值变号。随 \(p\) 光滑变化的 \(k\)-形式记作
\(\Omega^k(M)\)。对一形式 \(\alpha,\beta\)，
\[
(\alpha\wedge\beta)(v,w)
=\alpha(v)\beta(w)-\alpha(w)\beta(v).
\]
若 \(\alpha\) 为 \(k\)-形式、\(\beta\) 为 \(\ell\)-形式，
一般楔积由下面的交替化公式固定归一化：
\[
(\alpha\wedge\beta)(v_1,\ldots,v_{k+\ell})
=\frac1{k!\ell!}\sum_{\pi\in S_{k+\ell}}
\operatorname{sgn}(\pi)\,
\alpha(v_{\pi(1)},\ldots,v_{\pi(k)})
\beta(v_{\pi(k+1)},\ldots,v_{\pi(k+\ell)}).
\]
这里 \(S_{k+\ell}\) 是全部输入次序的排列，
\(\operatorname{sgn}(\pi)\) 在偶数次交换时为正、奇数次时为负。
交换两组输入需要 \(k\ell\) 次换位，因而
\(\alpha\wedge\beta=(-1)^{k\ell}\beta\wedge\alpha\)。
\end{definition}

在坐标中，\(\dd x^1\wedge \dd x^2(\partial_1,\partial_2)=1\)；
交换两个方向则为 \(-1\)。若 \(y=y(x)\)，将
\(\dd y^a=(\partial_i y^a)\dd x^i\) 代入楔积，得到二维情形
\[
\dd y^1\wedge \dd y^2
=\det\!\left(\frac{\partial(y^1,y^2)}
{\partial(x^1,x^2)}\right)\dd x^1\wedge \dd x^2.
\]
雅可比行列式由反对称性自动产生，这就是带方向面积积分的坐标变换律。

函数的微分已由 \(\dd f(v)=v(f)\) 定义。要对一形式继续求导，
就要求新的操作既与坐标变换相容，又保留
\(\dd(\dd f)=0\) 这一混合偏导交换的事实。

\begin{definition}[外微分]
在每张坐标图中，对
\(\alpha=\sum_I a_I\,\dd x^{i_1}\wedge\cdots\wedge \dd x^{i_k}\)，
定义
\[
\dd\alpha=\sum_I \dd a_I\wedge \dd x^{i_1}\wedge\cdots\wedge \dd x^{i_k}.
\]
下面验证这些局部公式拼成同一个
\(\dd:\Omega^k(M)\to\Omega^{k+1}(M)\)。
\end{definition}

上式为何与坐标无关？先在函数上，链式法则保证 \(\dd f\) 的定义不依赖图。
对新坐标 \(y^a\)，写
\(\dd y^a=(\partial_i y^a)\dd x^i\) 并按局部公式求导，得到
\(\dd(\dd y^a)=(\partial_j\partial_i y^a)\dd x^j\wedge \dd x^i=0\)：
二阶偏导对称，楔积反对称。局部公式还直接满足分级 Leibniz 法则
\[
\dd(\alpha\wedge\beta)
=\dd\alpha\wedge\beta+(-1)^k\alpha\wedge \dd\beta
\quad(\alpha\in\Omega^k(M)).
\]
于是把任意形式改写成 \(\dd y^a\) 的楔积后再求导，
与原坐标下的结果相同。相同的对称—反对称抵消也给出
\(\dd^2\alpha=0\)：系数的二阶偏导关于指标对称，
两个 \(\dd x\) 的楔积反对称，两者相乘后逐项相消。

例如平面上的一形式 \(\alpha=x\,\dd y-y\,\dd x\) 满足
\[
\dd\alpha=\dd x\wedge \dd y-\dd y\wedge \dd x=2\,\dd x\wedge \dd y.
\]
在单位圆盘上积分右边得 \(2\pi\)，沿边界逆时针计算左边也得
\(2\pi\)。这既是 Stokes 定理的一个直接检查，也说明边界方向
必须与面积方向配套。

\section{零旋度何时真的意味着存在势函数}

写一形式 \(\alpha=P(x,y)\,\dd x+Q(x,y)\,\dd y\)。
直接用外微分定义，
\(\dd\alpha=(\partial_xQ-\partial_yP)\dd x\wedge \dd y\)。
若 \(\alpha=\dd f\)，则 \(\dd\alpha=\dd^2f=0\)。
反方向却需要检查定义域的形状，而不是只检查偏导。

\begin{definition}[闭形式、恰当形式与星形区域]
满足 \(\dd\alpha=0\) 的形式称为闭形式；
若存在低一阶形式 \(\beta\) 使 \(\alpha=\dd\beta\)，
称 \(\alpha\) 为恰当形式。
平面区域 \(U\) 相对点 \(x_0\) 是星形的，若每个
\(x\in U\) 与 \(x_0\) 的整个线段都在 \(U\) 内。
\end{definition}

\begin{theorem}[一形式的局部势函数]
若 \(U\subset\RR^2\) 是星形开集，
\(P,Q\) 连续可微，且 \(\partial_xQ=\partial_yP\)，
则存在函数 \(f\) 使 \(\alpha=\dd f\)。
\end{theorem}
\begin{proof}
平移原点到星形中心，写 \(\alpha=\alpha_i(x)\dd x^i\)
（重复的 \(i=1,2\) 求和）并定义
\[
f(x)=\int_0^1 x^i\alpha_i(tx)\,\dd t.
\]
线段始终在 \(U\) 内，故积分有定义。
对第 \(j\) 个坐标求导，闭条件
\(\partial_j\alpha_i=\partial_i\alpha_j\) 给
\[
\partial_j f
=\int_0^1\left[\alpha_j(tx)
+t x^i\partial_j\alpha_i(tx)\right]\dd t
=\int_0^1\frac{\dd}{\dd t}\bigl[t\alpha_j(tx)\bigr]\dd t
=\alpha_j(x).
\]
所以 \(\dd f=\alpha\)。证明同时给出可实际计算的势函数，
而不只断言它存在。
\end{proof}

去掉星形条件就有反例。在穿孔平面
\(U=\RR^2\setminus\{0\}\) 上取
\[
\alpha=\frac{-y\,\dd x+x\,\dd y}{x^2+y^2}.
\]
它正是前一节极坐标中写出的局部 \(\dd\theta\)，
但右侧本身在整个穿孔平面有定义。
逐项求导得 \(\dd\alpha=0\)，所以它处处局部有势。
但沿单位圆 \((x,y)=(\cos t,\sin t)\) 有 \(\alpha=\dd t\)，
故 \(\oint\alpha=2\pi\)。若 \(\alpha=\dd f\) 且 \(f\) 是全局
单值函数，链式法则使闭路积分等于 \(f(\gamma(2\pi))-f(\gamma(0))=0\)，
矛盾。局部可积与全局单值是两个不同问题；
电磁势的规范和几何上的拓扑障碍都沿着这个区别发展。

\section{向量场的流与 李括号}
\label{sec:math-lie-bracket-flow}

切向量只属于一个点；向量场 \(X\) 给每个点分配一个切向量，
于是可沿其方向解常微分方程
\(\dot\gamma(t)=X(\gamma(t))\)。局部解的映射
\(\Phi_t(p)=\gamma_p(t)\) 称为 \(X\) 的流。
先沿 \(X\) 走一小段再沿 \(Y\) 走，与交换顺序一般不同。

\begin{definition}[向量场的李括号]
向量场 \(X,Y\) 都可作为函数上的微分算子。定义
\[
[X,Y](f)=X(Y(f))-Y(X(f)).
\]
这是另一个向量场，称李括号。
\end{definition}

为了核对它确为向量场，把 \(fg\) 代入上式并两次使用 Leibniz 法则；
含 \(X(f)Y(g)\) 与 \(Y(f)X(g)\) 的交叉项互相抵消，
剩下 \([X,Y](f)g+f[X,Y](g)\)。
由前一节“导子就是切向量”的定理，所得运算在每一点确是切向量。
若 \(X=X^i\partial_i,Y=Y^i\partial_i\)，对坐标函数作用得到
\[
[X,Y]=(X^j\partial_jY^i-Y^j\partial_jX^i)\partial_i.
\]
因此坐标向量场 \(\partial_i\) 两两对易。

\begin{definition}[标架与余标架]
开集上的局部标架是每一点的一组光滑切向量基
\(e_1,\ldots,e_n\)。满足 \(\theta^a(e_b)=\delta^a{}_b\)
的光滑一形式 \(\theta^1,\ldots,\theta^n\) 称为对偶余标架。
\end{definition}

标架可以不是坐标基。写
\([e_a,e_b]=C^c{}_{ab}e_c\)，
将外微分的一形式公式用于 \(e_a,e_b\)，得到
\[
\dd\theta^c(e_a,e_b)
=e_a[\theta^c(e_b)]-e_b[\theta^c(e_a)]
-\theta^c([e_a,e_b])
=-C^c{}_{ab}.
\]
因 \(C^c{}_{ab}=-C^c{}_{ba}\)，等价地
\[
\dd\theta^c=-\frac12C^c{}_{ab}\theta^a\wedge\theta^b.
\]
这条式子只测量标架是否来自坐标，尚未测量空间曲率。

在平面极坐标区域 \(r>0\) 取正交标架
\(e_r=\partial_r,\ e_\theta=r^{-1}\partial_\theta\)。
直接用上面的坐标表达算出
\([e_r,e_\theta]=-r^{-1}e_\theta\)。
其余标架是 \(\theta^r=\dd r,\theta^\theta=r\,\dd\theta\)，
所以 \(\dd\theta^\theta=\dd r\wedge \dd\theta
=r^{-1}\theta^r\wedge\theta^\theta\)。
平面仍是平的；非零括号源于这组正交基随位置转动。

\begin{exercise}
在 \(\RR^2\) 取 \(X=x\partial_y\)、\(Y=\partial_x\)。
从定义算 \([X,Y]\)，并说明为何两次沿流移动的顺序有影响。
\end{exercise}


球面上的位置和切向量已经能够不依赖单张坐标图来定义。
可是位于两点的切向量属于不同的空间，单靠坐标图无法定义它们之差；
这正是下一章引入联络的原因。
